\documentclass[10pt,a4paper]{amsart}
\usepackage[margin=19mm]{geometry}
\usepackage{amsmath,amssymb,mathtools}
\usepackage[scr=rsfs,cal=euler]{mathalfa}
\usepackage{xfrac}
\usepackage[unicode,hidelinks,bookmarks=false]{hyperref}
\newcommand{\ie}{i.e.\ }
\newcommand{\cf}{cf.\ }
\newcommand{\resp}{respectively\ }
\newcommand{\Subsection}{\subsection}
\providecommand{\nobreakdash}{\nobreak\hbox{-}}
\setlength{\emergencystretch}{2em}
\multlinegap0pt
\usepackage{bm}
\usepackage{bbm}
\usepackage{dsfont}
\usepackage{xspace}
\usepackage{cancel}
\usepackage{mathtools}
\usepackage{cleveref}
\usepackage{amsthm}
\makeatletter
\@ifundefined{theorem}{\newtheorem{theorem}{Theorem}}{}
\@ifundefined{definition}{\newtheorem{definition}[theorem]{Definition}}{}
\@ifundefined{lemma}{\newtheorem{lemma}[theorem]{Lemma}}{}
\makeatother
\def\Q{\mathcal{Q}}
\newcommand{\BRT}[0]{\mathsf{BRT}}
\newcommand{\NN}[0]{\mathbb{N}}
\newcommand{\nqvdash}{\operatorname{{?}{\nvdash}}}
\newcommand{\qvdash}{\operatorname{{?}{\vdash}}}
\title[Bounded Ramsey's theorem for triples in computability theory]{Bounded Ramsey's theorem for triples in computability theory}
\author{Ludovic Patey, Paul Shafer}
\date{Source: arXiv:2604.02092v1 (CC BY 4.0)}
\begin{document}
\maketitle

\noindent\emph{Source and licence.} Bounded Ramsey's theorem for triples in computability theory. Version arXiv:2604.02092v1 (CC BY 4.0), distributed under the Creative Commons Attribution 4.0 International licence, as recorded in the arXiv licence metadata for this exact version and shown on the abstract page https://arxiv.org/abs/2604.02092v1. Full rights notice: licenses/arxiv-2604.02092-CC-BY-4.0.txt.\par\smallskip

\noindent\textit{Editorial scope.} Counted unit: the Theorem whose source label is \texttt{None} in the article (source0.tex lines 458--461, source label \texttt{None}), delivered together with the article's own complete original proof of it (source0.tex lines 462--550, one \texttt{proof} environment of 89 lines). No mathematics is written, repaired or re-derived: statement and proof are the article's own text line for line. Required context retained uncounted: none. Every definition, display and label of those lines is retained unchanged. Omitted from this package and not redistributed: 1--457, 551--1025. Nothing inside a retained excerpt is omitted or altered: every retained body is delivered exactly as the article's lines are written, so no omission receipt of any kind accompanies this package. Citations: the retained excerpts carry no citation command at all, so no bibliography entry of the article is transcribed and no reference is redistributed. Reference adaptation: none was needed and none was made; every reference inside the delivered unit resolves inside it, so no unresolved reference or citation is produced. Printed slips kept literal: no printed word, symbol, hypothesis or number of the counted statement or of its proof was altered. Layout and class: the article's own class is \texttt{article} with packages including graphicx, amsthm, amsfonts, amsmath, amssymb, amsrefs, hyperref, cleveref, inputenc, url, tikz, tkz-graph, authblk, xcolor; the wrapper is the lane's pinned \texttt{amsart} editorial wrapper at 10pt on A4, which restates the theorem-like environments the retained text needs and adds the article's own macro definitions that the retained text needs (\texttt{\textbackslash BRT}, \texttt{\textbackslash NN}, \texttt{\textbackslash Q}, \texttt{\textbackslash nqvdash}, \texttt{\textbackslash qvdash}), with extra wrapper packages (bm, bbm, dsfont, xspace, cancel, mathtools, cleveref). The delivered numbering is the unit's own. Assets: no retained span contains a figure environment, a graphics-inclusion command, a PDF-page inclusion, a table or a TikZ picture; no image file is copied into this package, the package declares no asset dependency and no image file is redistributed with it. Licence: version arXiv:2604.02092v1 of this article is distributed under the Creative Commons Attribution 4.0 International licence, as recorded in the arXiv licence metadata for this exact version and shown on the abstract page https://arxiv.org/abs/2604.02092v1. A full rights notice is delivered at licenses/arxiv-2604.02092-CC-BY-4.0.txt. Characters: every retained excerpt is pure ASCII, so the delivered engine, XeLaTeX with its default Latin Modern text font, reports no missing character.
\begin{theorem}\label[theorem]{thm:delta2-cohesive-brt22-weakly-low}
Let $f : [\NN]^2 \to 2$ be a $\Delta^0_2$ instance of $\BRT^2_2$ and $P$ be of PA degree over~$\emptyset'$.
There is an infinite set~$G$ with $G' \leq_T P$ on which $f$ is stable and the limit is aways~0: $\lim_{y \in G} f(x,y) = 0$ for all $x \in G$.
\end{theorem}

\begin{proof}
Fix~$f$, and let $k \in \NN$ be such that there is no $f$-homogeneous set for color~1 of size~$k$.
\bigskip

\noindent
\textbf{$f$-tree.}
An \emph{$f$-tree} is a non-empty tree $T \subseteq \NN^{<\NN}$ such that for every~$\sigma \in T$, 
$\sigma$ is an increasing sequence that is $f$-homogeneous for color~1.  That is, for every $i < |\sigma|-1$, $\sigma(i) < \sigma(i+1)$, and for every $i < j < |\sigma|$, $f(\sigma(i), \sigma(j)) = 1$. 
Given $\sigma \in T$, we write $F(T, \sigma)$ for the set $\{ x \in \NN : \sigma \cdot x \in T \}$. Note that $F(T, \sigma)$ is not necessarily $f$-homogeneous for color~0. Requiring $T$ to be non-empty is equivalent to saying that $T$ contains the node $\epsilon$ of length~0. Also note that since $f$ is an instance of $\BRT^2_{2,k}$, every $f$-tree has depth at most~$k$. We write $T^*$ for the sub-tree of~$T$ of depth at most~$k-1$, that is, $T^* = \{ \sigma \in T : |\sigma| < k-1\}$. If $T$ is an infinite $f$-tree, then there is a node $\sigma \in T^*$ such that $F(T, \sigma)$ is infinite. We also write $D(T)$ for the union of the range of every node, that is, $D(T) = \bigcup_{\sigma \in T} F(T, \sigma)$.
\bigskip

\noindent
\textbf{Notion of forcing.}
Consider the notion of forcing whose conditions are pairs $(T, X)$, where
\begin{itemize}
    \item[(1)] $T \subseteq \NN^{<\NN}$ is a finite $f$-tree;
    \item[(2)] $X \subseteq \NN$ is an infinite set such that $D(T) < X$;
    \item[(3)] $X$ is low.
\end{itemize}
A condition $q = (S, Y)$ \emph{extends} $p = (T, X)$ if $T \subseteq S$, $Y \subseteq X$, $D(S) \setminus D(T) \subseteq X$ and finally $f(x,y) = 0$ whenever $\sigma \in T^*$, $x \in F(T, \sigma)$ and $y \in F(S, \sigma) \setminus F(T, \sigma)$.
Think of a condition $p = (T, X)$ as simultaneous Mathias conditions $p^{[\sigma]} = (F(T, \sigma), X)$ for every~$\sigma \in T^*$. In particular, if $q = (S, Y)$ extends $p = (T, X)$, then for every~$\sigma \in T^*$, $q^{[\sigma]}$ Mathias extends $p^{[\sigma]}$. A \emph{code} of a condition $(T, X)$ is a pair $(T, e)$ such that $e$ is a lowness index for~$X$, that is, $\Phi_e^{\emptyset'} = X'$. 

Any decreasing sequence of conditions $p_0 \geq p_1 \geq \cdots$ with $p_s = (T_s, X_s)$ induces an $f$-tree $T = \bigcup_s T_s$, and every $\sigma \in T^*$ induces a set $G_\sigma = F(T, \sigma) = \bigcup_s F(T_s, \sigma)$. As mentioned, if $T$ is infinite, then there is some~$\sigma \in T^*$ such that $G_\sigma$ is infinite.
Given $s \in \NN$ and $\sigma \in T_s^*$, $F(T_s, \sigma)$ is not necessarily $f$-homogeneous for color~0, and so neither is $G_\sigma$, but by the definition of condition extension, for every $x \in F(T_s, \sigma)$ and $y \in G_\sigma \setminus F(T_s, \sigma)$, $f(x, y) = 0$. It follows that every element in~$G_\sigma$ will have limit color~0 in~$G_\sigma$.
\bigskip

\noindent
\textbf{Forcing $\Sigma^0_1$ and $\Pi^0_1$ formulas.}
Let $p = (T, X)$ be a condition, $\varphi(G)$ be a $\Sigma^0_1$-formula and $\sigma \in T^*$. We shall always assume that $\varphi(G)$ is in normal form and induces a monotone formula $\varphi(\rho)$ on strings such that $\varphi(G)$ holds if and only if there is some~$\rho \prec G$ such that $\varphi(\rho)$ holds.
\begin{itemize}
    \item[(1)] $p \Vdash \varphi(G_\sigma)$ if $\varphi(F(T, \sigma))$ holds.
    \item[(2)] $p \Vdash \neg \varphi(G_\sigma)$ if for every finite set $H \subseteq X$, $\varphi(F(T, \sigma) \cup H)$ does not hold.
\end{itemize}
We say that $p$ \emph{decides} $\varphi(G_\sigma)$ if $p \Vdash \varphi(G_\sigma)$ or $p \Vdash \neg\varphi(G_\sigma)$. Note that the forcing relation for $\Pi^0_1$-formulas is stronger than the semantic one and yields the following lemma.

\begin{lemma}
Let $p_0 \geq p_1 \geq \cdots$ be a decreasing sequence of conditions with $p_s = (T_s, X_s)$.
Let $T = \bigcup_s T_s$ and $\sigma \in T^*$. If for every $\Sigma^0_1$-formula~$\varphi(G)$ there is some~$s$ such that $p_s$ decides $\varphi(G_\sigma)$, then $G_\sigma$ is infinite.
\end{lemma}
\begin{proof}
For every~$n \in \NN$, we can define a $\Sigma^0_1$-formula $\varphi_n(G) \equiv \exists x > n, x \in G$. There is no condition $p$ such that $p \Vdash \neg \varphi_n(G_\sigma)$, so if some $p_s$ decides $\varphi_n(G_\sigma)$, then $p_s \Vdash \varphi(G_\sigma)$, and unfolding the definition, there is some~$x > n$ such that $x \in F(T_s, \sigma)$.
Since $G_\sigma = \bigcup_s F(T_s, \sigma)$, it follows that $G_\sigma$ is infinite.
\end{proof}
\bigskip

\noindent
\textbf{Forcing question.}
Fix an enumeration $\varphi_0, \varphi_1, \dots$ of all $\Sigma^0_1$-formulas.
Since each condition represents a tree of Mathias conditions, we define a disjunctive forcing question.

\begin{definition}
Let $p = (T, X)$ be a condition and $g : T^* \to \NN$ be a function.
Let $p \qvdash \bigvee_{\sigma \in T^*} \varphi_{g(\sigma)}(G_\sigma)$ hold if for every function $h : X \to T^*$,
there is some $h$-homogeneous set~$H \subseteq X$ for some color~$\sigma \in T^*$ such that $\varphi_{g(\sigma)}(F(T, \sigma) \cup H)$ holds.
\end{definition}

Note that by compactness, the forcing question is $\Sigma^0_1(X)$ uniformly in the parameters of the formula.

\begin{lemma}\label[lemma]{lem:delta2-cohesive-brt22-weakly-low-forcing-question-spec}
Let $p = (T, X)$ be a condition and $g : T^* \to \NN$ be a function.
\begin{itemize}
    \item[(i)] If $p \qvdash \bigvee_{\sigma \in T^*} \varphi_{g(\sigma)}(G_\sigma)$, then there is an extension $q \leq p$ and some~$\sigma \in T^*$ such that $q \Vdash \varphi_{g(\sigma)}(G_\sigma)$.
    \item[(ii)] If $p \nqvdash \bigvee_{\sigma \in T^*} \varphi_{g(\sigma)}(G_\sigma)$, then there is an extension $q \leq p$ and some~$\sigma \in T^*$ such that $q \Vdash \neg \varphi_{g(\sigma)}(G_\sigma)$.
\end{itemize}
Moreover, a code for~$q$ and the case in which we are can be $P$-computed uniformly in a code for~$p$.
\end{lemma}
\begin{proof}\ 
\begin{itemize}
    \item[(i)] Suppose $p \qvdash \bigvee_{\sigma \in T^*} \varphi_{g(\sigma)}(G_\sigma)$ holds. Then by compactness, there is a finite set $E \subseteq X$ such that for every function $h : E \to T^*$, there is some $h$-homogeneous set~$H \subseteq E$ for some color~$\sigma \in T^*$ such that $\varphi_{g(\sigma)}(F(T, \sigma) \cup H)$ holds. 
    Let $h : E \to T^*$ be such that $h(y)$ is a node of maximal length~$\sigma \in T^*$ for which $\sigma \cup \{y\}$ is $f$-homogeneous for color~1. 
    
    We claim that for every~$y \in E$ and $x \in F(T, h(y))$, $f(x, y) = 0$. Fix~$y$ and~$x$. By definition of an $f$-tree, $h(y) \cup \{x\}$ is $f$-homogeneous for color~1, and by choice of $h(y)$, so is $h(y) \cup \{y\}$. It follows that if $f(x, y) = 1$, then $h(y) \cup \{x, y\}$ is $f$-homogeneous for color~1. Since $f$ is an instance of $\BRT^2_{2,k}$, it has size less than~$k$, so $|h(y)| < k-2$, but then $\sigma = h(y) \cdot x \in T^*$ is such that $\sigma \cup \{y\}$ is $f$-homogeneous for color~1, contradicting the maximality of $h(y)$. Thus $f(x, y) = 0$.

    Let $H \subseteq E$ be an $h$-homogeneous set for some color~$\sigma \in T^*$ such that $\varphi_{g(\sigma)}(F(T, \sigma) \cup H)$ holds. Let $S = T \cup \{ \sigma \cdot y : y \in H \}$. Then $q = (S, X \setminus [0, \max H])$ is an extension of~$p$ such that $q \Vdash \varphi_{g(\sigma)}(G_\sigma)$. Moreover, a code for~$q$ can be $\emptyset'$-computed from a code of~$p$.
    
    \item[(ii)] Suppose $p \nqvdash \bigvee_{\sigma \in T^*} \varphi_{g(\sigma)}(G_\sigma)$. Let $\Q$ be the $\Pi^0_1(X)$ class of all $h : X \to T^*$ such that for every~$\sigma \in T^*$ and every~$h$-homogeneous set~$H \subseteq X$, $\varphi_{g(\sigma)}(F(T, \sigma) \cup H)$ does not hold. By assumption, $\Q \neq \emptyset$. By the uniform low basis theorem, we can $\emptyset'$-compute a lowness index of $h \oplus X$ for some~$h \in \Q$. Moreover, one can $P$-compute a lowness index of an infinite $h$-homogeneous subset~$Y \subseteq X$ for some color~$\sigma$. The condition $q = (T, Y)$ is an extension of~$p$ such that $q \Vdash \neg \varphi_{g(\sigma)}(G_\sigma)$.
\end{itemize}
\end{proof}
\bigskip

\noindent
\textbf{Construction.}
Given a condition $p = (T, X)$, we let $g_p : T^* \to \NN$ be such that for each~$\sigma \in T^*$, $g_p(\sigma)$ is the least~$e \in \NN$ such that $p$ does not decide $\varphi_e(G_\sigma)$. Using \Cref{lem:delta2-cohesive-brt22-weakly-low-forcing-question-spec}, build a $P$-computable sequence of codes corresponding to a decreasing sequence of conditions $p_0 \geq p_1 \geq \cdots$ with $p_s = (T_s, X_s)$, such that for every~$s \in \NN$, there is some~$\sigma_s \in T_s^*$ such that $p_{s+1}$ decides $\varphi_{g_{p_s}(\sigma_s)}(G_{\sigma_s})$. Let $T  = \bigcup_s T_s$. 

We claim that $T$ is infinite. Indeed, if $T$ is finite, then by the pigeonhole principle, there is some~$\sigma \in T^*$ such that $\sigma = \sigma_s$ for infinitely many~$s$. By the definition of $g_p$, every $\Sigma^0_1(G_\sigma)$-formula is decided.  Thus $G_\sigma$ is infinite, so $T$ is infinite. 

Let $\sigma \in T^*$ be such that $G_\sigma$ is infinite. Note that an extension increases the size of~$G_\sigma$ at stage~$s$ only when $\sigma = \sigma_s$, so there are infinitely many~$s$ such that $\sigma = \sigma_s$. Again by the definition of~$g_p$, every $\Sigma^0_1(G_\sigma)$-formula is decided, so $G_\sigma' \leq_T P$. By the definition of condition extension, every element of~$G_\sigma$ has limit~0 in $G_\sigma$, so $f$ is stable on~$G_\sigma$.
This completes the proof of \Cref{thm:delta2-cohesive-brt22-weakly-low}.
\end{proof}

\end{document}
